\section{Validation design}
\label{sec:validation}
\label{sec:4}

\subsection{Case selection}\label{sec:4.1}

\begin{enumerate}
\def\labelenumi{\arabic{enumi}.}
\tightlist
\item
  \textbf{Coincident overpasses.} For every GPM orbit whose swath centre line crosses a region box between 1 June and 20 September, we estimate when the INSAT-3DR scan reaches that latitude (Section~\ref{sec:2.1}). Crossings within 5 min of an INSAT scan are kept.
\item
  \textbf{Convection screening.} Each candidate GPM file is screened for convection before INSAT data are downloaded. The densest convective cluster is found by counting GPM convective footprints in 1° bins and taking all convective footprints within about 150 km of the densest bin. A case needs at least 100 convective footprints in this cluster.
\item
  \textbf{Files per case:} the GPM 2A Ku orbit, the INSAT L1C scan nearest in time with its L2B CTP file, and daytime L1C scans of the case day and the days before and after for navigation.
\end{enumerate}

Cases were selected in two rounds: 2022--2024 for all four regions, then 2018--2021 for South and Central India, where the first round had few usable cases. In total 63 cases were processed.

\subsection{Colocation}\label{sec:4.2}

Each case is processed on a 0.02° grid over its region box. INSAT data are corrected (Section~\ref{sec:3}) and also resampled uncorrected. GPM footprints are assigned to grid cells by nearest neighbour within 3.5 km; storm tops use their slanted-beam positions (Section~\ref{sec:2.2}). Only cells inside the GPM swath are compared. Analysis is limited to within 150 km and 250 km of the centre of the densest convective cluster.

\subsection{Offset measures}\label{sec:4.3}

Because the correction is a displacement, we measure displacements directly, using two independent methods.

\textbf{Field correlation.} We find the shift of the INSAT field that best correlates with the GPM field:

\begin{itemize}
\tightlist
\item
  primary pair: \(-\)BT against GPM storm-top height (both describe storm tops);
\item
  secondary pair: \(-\)BT against near-surface reflectivity.
\end{itemize}

Integer shifts up to ±10 cells (±20 km) are tested, and the best one is refined to a fraction of a cell by fitting a parabola through the correlation at the peak and its two neighbours, separately along each axis.

\textbf{Cell cores.} Each GPM deep-echo cell (connected cells with storm top \textgreater{} 10 km, at least 8 grid cells) is paired with the INSAT cold core within 25 km of its centroid. The cold core is the centroid of pixels within 3 K of the local minimum BT. The case value is the median over its cells. This measure does not depend on the shape or extent of the anvil.

\textbf{Along and across the line of sight.} Each offset is split into a component along the line of sight (positive when INSAT must move towards the satellite to match GPM) and a component across it, using the VAA at the cluster centre. Parallax correction should change only the along-look component. The cross-look component is a control.

\textbf{Overlap scores.} We also report:

\begin{itemize}
\tightlist
\item
  intersection over union (IoU) between INSAT deep cores (BT \textless{} 220 K) and GPM deep echoes (storm top \textgreater{} 8 km);
\item
  Pearson correlation and mutual information between cold cloud (BT \textless{} 235 K) and echo (reflectivity \textgreater{} 25 dBZ).
\end{itemize}

\subsection{Usable cases}\label{sec:4.4}

A case is usable if its convective cluster has at least 100 footprints and at least one field and window gives a reliable match. A match is reliable when the best shift is inside the search range (below 9.5 cells) both before and after correction. Matches at the edge of the search usually mean the two fields describe different parts of the storm.

\subsection{Experiments}\label{sec:4.5}

Each case is run in four states:

\begin{longtable}[]{@{}lll@{}}
\caption{Correction states.}\label{tab:S}\\
\toprule\noalign{}
State & Navigation correction & Parallax correction \\*
\midrule\noalign{}
\endfirsthead
\bottomrule\noalign{}
\endlastfoot
Raw & no & no \\*
Navigation only & per-year monsoon correction (Section~\ref{sec:3.6}) & no \\*
Parallax only & no & yes \\*
Navigation + parallax & per-year monsoon correction & yes \\*
\end{longtable}

Two further tests address the residual offset:

\begin{itemize}
\tightlist
\item
  \textbf{Navigation from GPM.} If the residual cell offset after parallax correction were a navigation error, a navigation correction fitted to it should improve both offset measures. For each case we take the median cell offset (converted to east and north) of all \emph{other} cases and apply it with parallax correction (leave one out, so no case is corrected with its own data).
\item
  \textbf{Synthetic storm.} A storm with known height and position is displaced by its computed parallax and put through the same gridding and matching code, to check that the method recovers the full parallax and zero after correction.
\end{itemize}

To separate sources of residual offset, we also compare INSAT cold-core heights with GPM maximum echo-top heights in each matched cell, and repeat the cell analysis with GPM storm tops at their surface footprints instead of on the beam.

\subsection{Sensitivity}\label{sec:4.6}

An error \(\delta h\) in cloud-top height shifts the corrected position by about \(\delta h\tan(\mathrm{VZA})\): about 0.4 km per km of height error at Mumbai and 0.7 km at Agartala. We quantify this with Monte Carlo perturbations of cloud-top height (standard deviation 1 km, 100 members) and with fixed biases of ±1 and ±2 km.

\subsection{Statistics}\label{sec:4.7}

Uncertainties are 95 \% percentile bootstrap intervals from 10 000 resamples. For the field measure, cases are resampled; for the cell measure, cases are resampled with all their cells (a cluster bootstrap), because cells of one storm are not independent. Paired comparisons between correction states use the absolute offset of each case: the median change with its bootstrap interval, the Wilcoxon signed-rank test and the exact sign test on the number of cases that improve. For the cell measure the paired test uses the per-case medians. The dependence of the raw offset on viewing angle is tested with the Spearman rank correlation and the least-squares slope against tan(VZA), with a bootstrap interval.

